% Part I > B. Open Societal Problems: the problems the clusters name, each with
% its data. Integrity (the Poison Fountain, unchanged), credit and careers,
% understanding. Access and funding are held for the closing's open problems.
% talk-v2, 2026-11-11.
\section{Open Societal Problems}

% [1.5 min] The fact of rampant AI use, and what it does to checking. Left:
% monthly share of arXiv math preprints that (blue) acknowledge AI use (Epoch),
% (orange) read as AI-written without disclosing AI writing (Martin; series
% starts Aug 2025), (purple) acknowledge AI for formalization (Epoch).
% Right: Kyle's four points, each with its numbers.
% Kyle's original bullets:
%   - Rampant AI use has led to exponentially increasing submissions to journals
%     and preprint servers, with an increasing number disclosing AI use
%   - However, many do not disclose. An increasing fraction are even thought to
%     be generated in whole by AI.
%   - OpenAI set the standard of accompanying results with formal certificates,
%     but relatively few are doing this, in particular in machine learning.
%   - (Other significant statistics regarding published work using AI)
% Data: research/03-publishing-crisis.md, "AI use, disclosure and formalization".
% Colours: tlBlue / tlOrange / tlPurple pass the dataviz palette validator.
\begin{frame}{The Poison Fountain}
  \centering
  \begin{tikzpicture}[x=1cm, y=1cm,
      claimnum/.style={circle, fill=PennRed, text=white, inner sep=0pt, minimum size=0.42cm,
                       font=\sffamily\scriptsize\bfseries},
      claim/.style={anchor=north west, text width=4.95cm, align=flush left, inner sep=0pt,
                    font=\sffamily\scriptsize, text=black!75}]
    \useasboundingbox (-0.4,0.15) rectangle (14.21,-6.35);
    % ---- chart -------------------------------------------------------------
    \node[anchor=north west, inner sep=0pt, font=\sffamily\scriptsize\bfseries, text=black!70]
      at (0.0,0.1) {Share of arXiv math preprints, by month};
    \begin{axis}[at={(1.0,-0.75)}, anchor=north west, scale only axis, width=6.55cm, height=4.75cm,
        xmin=-0.5, xmax=21.5, ymin=0, ymax=52, clip=false,
        xtick={0,6,12,18,21}, xticklabels={Jan '25, Jul '25, Jan '26, Jul '26, Oct},
        ytick={0,10,20,30,40,50}, yticklabel={\pgfmathprintnumber[assume math mode=true]{\tick}\%},
        tick label style={font=\sffamily\tiny, text=black!60},
        axis line style={draw=black!30}, tick style={draw=black!30},
        ymajorgrids, grid style={draw=black!8},
        axis x line*=bottom, axis y line*=left,
        legend style={at={(0.02,0.98)}, anchor=north west, draw=none, fill=white, fill opacity=0.9,
                      text opacity=1, font=\sffamily\tiny, row sep=-1pt, inner sep=1pt},
        legend cell align=left,
        every axis plot/.append style={line width=1.1pt, line join=round, mark size=0.75pt}]
      % disclosed AI use, any purpose (Epoch AI, recomputed from the 9 Nov 2026 CSV)
      \addplot[color=tlBlue, mark=*] coordinates {
        (0,0.35) (1,0.49) (2,0.23) (3,0.77) (4,0.80) (5,0.92) (6,0.86) (7,1.12) (8,1.74) (9,1.50) (10,1.36)
        (11,1.91) (12,1.93) (13,2.39) (14,3.68) (15,3.89) (16,6.64) (17,8.43) (18,15.42) (19,26.33) (20,38.55) (21,48.05)};
      \addlegendentry{Acknowledge AI use}
      % apparent AI writing with no disclosure of AI writing (Martin, 3 Nov 2026)
      \addplot[color=tlOrange, mark=*] coordinates {
        (7,4.9) (8,4.2) (9,5.0) (10,6.3) (11,5.5) (12,7.4) (13,7.4) (14,9.7) (15,12.2) (16,14.4)
        (17,16.2) (18,19.1) (19,21.1) (20,20.6) (21,17.5)};
      \addlegendentry{Read as AI-written, not disclosed}
      % AI acknowledged for formalization (Epoch AI)
      \addplot[color=tlPurple, mark=*] coordinates {
        (0,0) (1,0) (2,0) (3,0.03) (4,0) (5,0) (6,0) (7,0.03) (8,0.03) (9,0.02) (10,0.05)
        (11,0.05) (12,0.06) (13,0.17) (14,0.14) (15,0.39) (16,0.16) (17,0.55) (18,0.68) (19,1.31) (20,1.81) (21,1.88)};
      \addlegendentry{AI used for formalization}
      % end values, in text ink beside each line
      \node[anchor=west, font=\sffamily\tiny\bfseries, text=black!75] at (axis cs:21.4,48.05) {48\%};
      \node[anchor=west, font=\sffamily\tiny\bfseries, text=black!75] at (axis cs:21.4,17.5) {17.5\%};
      \node[anchor=west, font=\sffamily\tiny\bfseries, text=black!75] at (axis cs:21.4,1.88) {1.9\%};
    \end{axis}
    % ---- claims ------------------------------------------------------------
    \foreach \n/\y in {1/-0.05, 2/-1.65, 3/-3.25, 4/-4.85} \node[claimnum] at (8.95,\y-0.17) {\n};
    \node[claim] at (9.35,-0.05)
      {\textbf{\color{black}Use is exploding.} Disclosed AI use went from 0.4\% to 48\% of math
       preprints in 21 months, doubling about every three months. Math submissions are up
       75\% on a year earlier (Sep 2026).};
    \node[claim] at (9.35,-1.65)
      {\textbf{\color{black}Much of it is undisclosed.} About 1 in 6 math preprints read as AI-written
       without saying so (Martin, Oct 2026). At ICLR 2026, 9\% of submissions were mostly AI-written.};
    \node[claim] at (9.35,-3.25)
      {\textbf{\color{black}Few certify.} OpenAI shipped Lean proofs with 42\% of its 719 manuscripts.
       Across arXiv math, 1.9\% acknowledge AI for formalization; in stat.ML, 0.4\% even mention Lean.};
    \node[claim] at (9.35,-4.85)
      {\textbf{\color{black}Review is overwhelmed.} 21\% of ICLR 2026 reviews were fully AI-generated.
       arXiv capped submissions (1 Oct); six probability journals paused (19~Oct\,--\,8~Nov).};
  \end{tikzpicture}
  \slidesources{Epoch AI (CC BY 4.0, 9 Nov 2026) \textbullet{} K.~Martin, \emph{AI use in arXiv mathematics papers}, kimballmartin.github.io/ai-use (3 Nov 2026; 17.5\% of Oct 2026 math preprints) \textbullet{} our arXiv count \textbullet{} Pangram, ICLR 2026 (self-reported) \textbullet{} AcadMathSci, Oct 2026}
\end{frame}

% [1.5 min] Credit and careers: the "AI capable, negative" cluster's words
% (scooped, junior mathematicians, hiring committees, disclosures). Rows 1-2
% from the former backup slide "Credit, disclosure and confidentiality"; the
% disclosure row is on the Poison Fountain. Paper mills: Dolores-Cuenca, corpus
% post 1043. Careers: anonymous PDE postdoc, corpus post 1893; hiring
% responses from AcadMathSci s3.2, s4.4.
\begin{frame}{Credit and careers}
  \footnotesize\sffamily
  \renewcommand{\arraystretch}{1.35}
  \begin{tabularx}{\linewidth}{@{}p{2.1cm}YYp{1.5cm}@{}}
    \toprule
    \tabhead{Problem} & \tabhead{Case} & \tabhead{Response} & \tabhead{Status}\\
    \midrule
    \rowhead{Credit} & AI reuses known ideas without attribution: 8 of Aletheia's 13 meaningful Erd\H{o}s results were already in the literature. Kolda: ``Human--AI scrapple''
      & Leiden Declaration (Jun 2026, endorsed by the IMU): human responsibility, attribution & \stpartial\\
    \rowhead{Priority} & Navier--Stokes priority dispute (C\'ordoba \& Mart\'inez-Zoroa; Buckmaster \& Alp\"oge). A result named in a talk title was pre-empted by a hasty AI-generated paper (anecdotal)
      & AGMAI release guidelines (29 Sep): review before announcement, release prompts. Labs withheld prompts & \stopen\\
    \rowhead{Paper mills} & A company could solve open questions with AI and sell the authorship of correct papers
      & None yet & \stopen\\
    \rowhead{Careers} & \textbf{``I no longer trust the profession to offer a fair and intelligible path forward''} (anonymous PDE postdoc)
      & Judge a few best outputs; viva-style questioning for hiring (AcadMathSci) & \stpartial\\
    \bottomrule
  \end{tabularx}
  \slidesources{Reisinger et al., AcadMathSci, Oct 2026 (\S3.2, \S3.4, \S4.4, A.3, C.1) \textbullet{} Kolda, Silicon Reckoner, Feb 2026 \textbullet{} Proofs \& Prompts posts 1043 (Dolores-Cuenca) and 1893 (anonymous)}
\end{frame}

% [1 min] Understanding: the "cognitive", "loss" and "obsolete" words. Rows 1-2
% from the former backup slide "Training, access and narrowing" (its access
% row is held for the closing). Row 3: Hairer (NYT interview, Feb 2026,
% secondary wording, research/05) and Moshkovitz on one OpenAI paper
% (research/03 s3).
\begin{frame}{Understanding}
  \footnotesize\sffamily
  \renewcommand{\arraystretch}{1.35}
  \begin{tabularx}{\linewidth}{@{}p{2.1cm}YYp{1.5cm}@{}}
    \toprule
    \tabhead{Problem} & \tabhead{Case} & \tabhead{Response} & \tabhead{Status}\\
    \midrule
    \rowhead{Deskilling} & Students given unrestricted GPT-4 scored 17\% lower once access was removed; a tutor with guardrails largely avoided the harm (Bastani et al.)
      & Vivas and in-person exams; AI-permitted work graded on formulation and validation & \stpartial\\
    \rowhead{Narrowing} & AI adoption shrinks topic range by 4.6\% and engagement by 22\% (Hao et al., 41M papers). Agendas drift toward ``problems easily amenable to current AI''
      & Reward understanding and long-term contribution over volume (AcadMathSci) & \stopen\\
    \rowhead{Opaque output} & Hairer: AI proofs detail the easy steps and skim the hard ones, and are hard to trust. Moshkovitz, on one of OpenAI's papers: impossible to read without AI help
      & Leiden: preserve human understanding. Proposals to publish the certificate apart from the explanation & \stopen\\
    \bottomrule
  \end{tabularx}
  \slidesources{Bastani et al., PNAS \textbullet{} Hao et al., arXiv 2412.07727 \textbullet{} Reisinger et al., AcadMathSci, Oct 2026 (\S3.2, \S3.3) \textbullet{} Hairer, NYT interview, Feb 2026 (secondary) \textbullet{} Leiden Declaration, Jun 2026}
\end{frame}

% [1 min] Bridge (Kyle, 2026-11-11): most of these problems are ongoing and have
% no clear solution, which puts the spotlight on the one benefit the essays
% agree on. Statuses for the three integrity items are a judgement from the
% Poison Fountain evidence (caps and pauses are temporary; disclosure rules are
% widely ignored; review capacity is unchanged); the rest are copied from the
% two tables above. Corpus count: posts in the 152-post analysis set whose text
% matches formaliz*/formalis*, "proof assistant" or the word Lean (case-
% sensitive): 24/66 capable-positive, 10/37 capable-negative, 11/38 doubted,
% 2/11 capability not addressed; 47/152 in all. Distinctive words: the
% capable-positive word cloud (data/wordcloud_terms.csv).
\begin{frame}{Ongoing, without clear solutions}
  \centering
  \begin{tikzpicture}[x=1cm, y=1cm,
      phead/.style={anchor=north west, inner sep=0pt, font=\sffamily\small\bfseries, text=black!80},
      plist/.style={anchor=north west, text width=4.3cm, align=flush left, inner sep=0pt,
                    font=\sffamily\footnotesize, text=black!85}]
    \useasboundingbox (0,0) rectangle (14.3,-5.9);
    \foreach \x/\h in {0/Integrity, 4.85/Credit and careers, 9.7/Understanding} {
      \fill[black!20] (\x,0) rectangle ++(4.45,-0.05);
      \node[phead] at (\x,-0.18) {\h};
    }
    \node[plist] at (0,-0.68) {%
      Volume of submissions\hfill\stpartial\par\smallskip
      Undisclosed AI\hfill\stpartial\par\smallskip
      Overwhelmed review\hfill\stopen};
    \node[plist] at (4.85,-0.68) {%
      Credit\hfill\stpartial\par\smallskip
      Priority\hfill\stopen\par\smallskip
      Paper mills\hfill\stopen\par\smallskip
      Careers\hfill\stpartial};
    \node[plist] at (9.7,-0.68) {%
      Deskilling\hfill\stpartial\par\smallskip
      Narrowing\hfill\stopen\par\smallskip
      Opaque output\hfill\stopen};
    \node[anchor=north, inner sep=0pt, font=\sffamily\small] at (7.15,-2.7)
      {None of these ten problems has a solution in place.};
    % ---- the one benefit ------------------------------------------------------
    \draw[black!15, line width=0.4pt] (0,-3.3) -- (14.3,-3.3);
    \node[anchor=north west, inner sep=0pt, font=\sffamily\small\bfseries, text=PennBlue] at (0,-3.5)
      {The one benefit that comes up across the board: formalization};
    \foreach \x/\col/\pct/\name in {0/clPos/36\%/{AI capable, positive}, 4.85/clNeg/27\%/{AI capable, negative}, 9.7/clDoubt/29\%/{AI doubted, negative}} {
      \node[anchor=north west, inner sep=0pt, font=\sffamily\bfseries\fontsize{18}{20}\selectfont, text=\col] at (\x,-4.0) {\pct};
      \node[anchor=north west, inner sep=0pt, text width=2.9cm, font=\sffamily\scriptsize, text=black!70] at (\x+1.35,-4.05)
        {of the essays in \textbf{\color{\col}\name} discuss it};
    }
    \node[anchor=north west, text width=14.2cm, align=flush left, inner sep=0pt, font=\sffamily\scriptsize, text=black!60] at (0,-5.0)
      {47 of 152 essays mention formalization, Lean or proof assistants. In the most positive cluster,
       \emph{checking}, \emph{formalization}, \emph{errors} and \emph{verify} are among the most distinctive words.};
  \end{tikzpicture}
  \slidesources{The three slides above \textbullet{} Proofs \& Prompts corpus: keyword count over the 152-post analysis set; word statistics from cluster.py}
\end{frame}
